\section{Introduction}
\label{sec:introduction}

Software in regulated domains such as finance, automotive, and power must
conform to extensive regulations, standards, and internal
policies~\cite{badrzadeh2015challenges,tabani2019assessing,xue2024llm4fin,zhang2025comprehensive}.
Regulatory compliance testing provides a final opportunity to detect
violations before deployment~\cite{sommerville2005ingenieria}. Yet producing
these tests remains largely manual: experts must interpret regulatory
documents, formalize implicit conditions, and connect test cases to
system interfaces~\cite{castellanos2022compliance,jain2025complexity}.
Regulations also contain cross-rule dependencies and context-sensitive
terminology rather than isolated requirements. A usable workflow must therefore
preserve source provenance, represent pre- and post-operation states, and
compose related rules into structured test scenarios.

Model-based approaches derive tests from domain models, graphs, or formal
representations and offer interpretability and controllability, but adapting
their artifacts can require substantial engineering
effort~\cite{wang2022automatic,stefani2025automated,zyberaj2025genai,yang2025llmcfg}.
LLM-based approaches automate more of the process through prompting,
retrieval-augmented generation, or multi-agent
collaboration~\cite{lewis2020retrieval,castellanos2022compliance,boukhlif2024towards,korraprolu2025test,hasan2025automatic,masuda2026generating},
but can produce hallucinated or inconsistent outputs. Hybrid methods guide
LLMs with ontologies, schemas, or formal specifications, improving structure
while still depending on manually constructed domain
knowledge~\cite{xue2024llm4fin,chen2024automated,sun2025compliance,shrestha2025less}.

The bottleneck therefore shifts from test generation to reusable knowledge
construction. The challenge is not simply a lack of domain facts. Direct
generation collapses rule interpretation, representation choice, and test
design into one opaque output, making errors difficult to localize and
improvements difficult to reuse. At the other extreme, manually engineered
artifacts improve control but must be rebuilt and maintained for each domain.
We therefore need an explicit boundary between authoritative rule semantics and
reusable generation guidance. This raises two questions: which intermediate representation is
structured enough for downstream testing, and how can its supporting domain
knowledge be obtained without recreating the same manual burden? We use formal
requirements because they balance expressiveness, machine interpretability,
and testability~\cite{gropler2021nlp,yang2025requirements}.

We address this bottleneck with \ourtool, a framework for
\underline{R}equirements \underline{A}uto-\underline{F}ormalization and
Compliance \underline{T}est Generation. \ourtool\ separates two knowledge
layers. The input regulation remains authoritative for rule-level obligations,
thresholds, and exceptions. Implicit regulatory knowledge in LLMs is explicated
into three reusable meta-level artifacts: a domain meta-model, a requirements
representation, and testability constraints. A small construction
corpus of documents and historical tests is optional and supplies evidence
about scope, terminology, granularity, and output format. \ourtool\ injects
these artifacts at the corresponding stages to formalize regulations, assess
testability, and generate compliance test cases. The meta-model
defines domain concepts, relations, operations, and state fields; the
requirements representation defines a constrained \textit{KEY-OP-VALUE}
grammar; and the testability artifact defines OCL-style decision predicates.
Generated formal requirements retain source-rule evidence and explicitly
encode their pre- and post-operation states. Knowledge construction is
completed and frozen before evaluation documents enter the generation phase,
preventing held-out rule content from leaking into the reusable artifacts.
This layered design makes errors attributable to vocabulary, formalization,
testability, or generation and permits artifact inspection and correction
without changing the authoritative rule text.

\ourtool\ constructs candidates with multiple LLMs and applies an \strategy\
strategy. Purification retains majority-supported meta-model elements, whereas
aggregation preserves complementary representation and testability knowledge.
The assignment follows artifact-specific failure costs: unsupported meta-model
elements can corrupt downstream stages, while deleting valid representation or
testability knowledge reduces coverage. Artifacts and provenance are retained
for inspection.

We evaluate \ourtool\ on 11 financial, road-traffic, and power datasets. The
controlled E2E baseline matches backends and documents but omits explicit
artifacts. On six financial datasets, \ourtool\ achieves 91.7\% F1/87.1\% BSC,
23.6 F1 points above the strongest automated baseline and higher on every
dataset; HITL reaches 96.2\% F1. Automated \ourtool, HITL, and experts require
2.53 elapsed hours, 5.03 elapsed hours (7.16 person-hours), and 42.56
person-hours. Secondary analyses show 90.2\% corpus-free F1, gains from explicit
artifacts and operator assignment, wins on 10 of 11 backbones, and
re-instantiated F1 of 92.7\%/87.2\% in road traffic/power.

Beyond the metric gains, \ourtool\ turns transient model outputs into reusable
engineering assets. Explicit artifact contracts make knowledge inspectable and
traceable, support targeted revision before generation, and provide stable
interfaces across changing regulations, models, and domains. This design
supports repeated compliance workflows and incremental maintenance. Because the
artifacts are decoupled from individual rules, teams can reuse them across
document revisions and compare outputs across LLM backbones. Their shared
structure also gives engineers and domain experts a common review surface.

Together, these results show the practical value of reusable, inspectable
artifacts for compliance test generation.

The contributions of this paper are:
\begin{enumerate}[leftmargin=*,align=left,label=(\arabic*), itemsep=0pt, parsep=0pt, nosep]
    \item \textbf{Automated Compliance-Test Framework.}
    \ourtool\ separates authoritative rule semantics from reusable artifacts
    for traceable formalization, testability assessment, and test generation.

    \item \textbf{Risk-Adaptive Knowledge Explication.}
    Our \strategy\ strategy assigns purification or aggregation by artifact type
    to consolidate three complementary artifacts from multiple LLM outputs.

    \item \textbf{Cross-Domain Evaluation.}
    We evaluate three domains with ablations, expert review, step-level and
    backbone analyses, and a shared cost, time, and run-variation protocol.
\end{enumerate}
